\section{Weighted phase banks and conjugation conventions}
\label{sec:tpc40-weighted-banks}

Fix positive integers \(m\) and \(R\).  A phase bank is a matrix
\begin{equation}
 \Phi=(\phi_{j\alpha})_{1\le j\le m,\,1\le\alpha\le R}
 \in\C^{m\times R},
 \qquad |\phi_{j\alpha}|=1.
 \label{eq:tpc40-phase-bank}
\end{equation}
Its rows are measurements and its columns are the \(R\) signatures to be
recovered with the same unit gain.  Let
\begin{equation}
 w_j>0,
 \qquad \sum_{j=1}^{m}w_j=1,
 \qquad W:=\operatorname{diag}(w_1,\ldots,w_m).
 \label{eq:tpc40-row-weights}
\end{equation}

We use the following complex conventions throughout.  For a column \(x\),
\(x^T\) denotes ordinary transpose, \(x^*=\overline{x}^{\,T}\) denotes
conjugate transpose, and
\begin{equation}
 \langle x,y\rangle=x^*y
 \label{eq:tpc40-inner-product-convention}
\end{equation}
is conjugate linear in the first variable and linear in the second.  The
recovery coefficients form a column \(a\in\C^m\), but the recovery rule uses
ordinary transpose:
\begin{equation}
 a^T\Phi=\one_R^T.
 \label{eq:tpc40-exact-recovery}
\end{equation}
Thus no conjugate is silently inserted into the measurement rule.  Taking
the conjugate transpose of \eqref{eq:tpc40-exact-recovery} gives the
equivalent equation
\begin{equation}
 \Phi^*\overline a=\one_R.
 \label{eq:tpc40-exact-recovery-adjoint}
\end{equation}

Define the weighted normalized amplification and weighted Gram matrix by
\begin{align}
 A_W(a)
 &:=a^*W^{-1}a
   =\sum_{j=1}^{m}\frac{|a_j|^2}{w_j},
 \label{eq:tpc40-weighted-amplification}\\
 G_W(\Phi)
 &:=\Phi^*W\Phi.
 \label{eq:tpc40-weighted-Gram}
\end{align}
When \(w_j=1/m\), these reduce to
\begin{equation}
 A_W(a)=m\|a\|_2^2,
 \qquad
 G_W(\Phi)=\frac1m\Phi^*\Phi.
 \label{eq:tpc40-uniform-specialization}
\end{equation}
We normally abbreviate them to \(A\) and \(G\).  The physical-interface
sections use the typographic form \(\cA\) for the same amplification.

It is useful to absorb the weights into
\begin{equation}
 D:=W^{1/2}\Phi,
 \qquad
 b:=W^{-1/2}a,
 \qquad
 z:=\overline b=W^{-1/2}\overline a.
 \label{eq:tpc40-weighted-normalization}
\end{equation}
Then
\begin{equation}
 G=D^*D,
 \qquad
 A=\|b\|_2^2=\|z\|_2^2,
 \qquad
 D^*z=\one_R.
 \label{eq:tpc40-normalized-recovery-system}
\end{equation}
The last identity is the adjoint form of \(b^TD=\one_R^T\).

Because every phase has modulus one and the weights sum to one,
\begin{equation}
 G\succeq0,
 \qquad G_{\alpha\alpha}=1,
 \qquad \operatorname{tr}G=R.
 \label{eq:tpc40-Gram-basic-facts}
\end{equation}
In particular,
\begin{equation}
 \lambda_{\max}(G)\ge1,
 \label{eq:tpc40-Gram-trace-lower-bound}
\end{equation}
with equality if and only if \(G=I_R\).  The latter assertion follows
because all \(R\) nonnegative eigenvalues must then equal their common
upper bound \(1\).
